% Zone „Das kennst du schon" – Quadratische Gleichungen, Kl. 9 Oberschule
\hypertarget{zone}{}\einheitenkopf{Das kennst du schon}
\renewcommand{\theaufgabe}{Z\arabic{aufgabe}}\setcounter{aufgabe}{0}

\begin{aufgabe}{Ich kann Zahlen quadrieren, auch negative Zahlen und Dezimalzahlen. Berechne.}
\begin{teilezwei}
\tz $5^2 =$ \leerfeld & \tz $9^2 =$ \leerfeld \\
\tz $(-8)^2 =$ \leerfeld & \tz $0{,}4^2 =$ \leerfeld \\
\tz $2{,}5^2 =$ \leerfeld & \tz $3 \cdot (-2)^2 =$ \leerfeld \\
\tz $10 - 4^2 =$ \leerfeld & \\
\end{teilezwei}
\end{aufgabe}

\begin{aufgabe}{Ich kann Klammern zuerst rechnen und Punkt vor Strich beachten, auch mit negativen Zahlen. Berechne.}
\begin{teilezwei}
\tz $3 \cdot (2 + 4) =$ \leerfeld & \tz $(7 - 5) \cdot 6 =$ \leerfeld \\
\tz $(-3) \cdot (-3 + 8) =$ \leerfeld & \tz $(-4) \cdot (-4 - 2) =$ \leerfeld \\
\tz $(-1 - 3) \cdot (-1 + 2) =$ \leerfeld & \tz $(-3)^2 + 5 \cdot (-3) + 6 =$ \leerfeld \\
\end{teilezwei}
\end{aufgabe}

\begin{aufgabe}{Ich kann eine lineare Gleichung durch Umformen lösen. Löse die Gleichung. Schreibe hinter den Strich, was du auf beiden Seiten tust.}
\begin{gleichungsraster}[2]
\gl{x - 6 = 3} & \gl{4x = 28} \\ \rule{0pt}{2mm} & \\
\gl{-x = 8} & \gl{x + 2 = -5} \\ \rule{0pt}{2mm} & \\
\gl[3]{2x + 7 = -3} & \gl[3]{5x - 3 = 2x + 9} \\ \rule{0pt}{2mm} & \\
\end{gleichungsraster}
\end{aufgabe}

\begin{aufgabe}{Ich kann durch Einsetzen prüfen, ob eine Zahl eine Gleichung löst. Setze die Zahl ein, rechne beide Seiten aus und kreuze an.}
\begin{geruest}
\gz{a}{$x = 3$ \ in \ $2x + 1 = 7$}{\kreuz{(wA)}\kreuz{(fA)}}
\gz{b}{$x = 5$ \ in \ $x - 2 = 4$}{\kreuz{(wA)}\kreuz{(fA)}}
\gz{c}{$x = -2$ \ in \ $3x + 10 = 4$}{\kreuz{(wA)}\kreuz{(fA)}}
\gz{d}{$x = -3$ \ in \ $4 - 2x = 10$}{\kreuz{(wA)}\kreuz{(fA)}}
\gz{e}{$x = -1$ \ in \ $5x + 2 = 3x - 4$}{\kreuz{(wA)}\kreuz{(fA)}}
\end{geruest}
\end{aufgabe}

\begin{aufgabe}{Ich kann Quadratwurzeln ziehen und Näherungswerte runden. Berechne. Geht die Wurzel nicht auf, runde auf zwei Stellen nach dem Komma. Gibt es keine Wurzel, schreibe „gibt es nicht“.}
\begin{teilezwei}
\tz $\sqrt{81} =$ \leerfeld & \tz $\sqrt{144} =$ \leerfeld \\
\tz $\sqrt{0{,}25} =$ \leerfeld & \tz $\sqrt{7} \approx$ \leerfeld \\
\tz $\sqrt{0} =$ \leerfeld & \tz $\sqrt{-16} =$ \leerfeld \\
\end{teilezwei}
\end{aufgabe}

\begin{aufgabe}{Ich kann an der Scheitelpunktform den Scheitel, die Öffnung und die Zahl der Nullstellen ablesen. Gib den Scheitel an, kreuze an, wie die Parabel geöffnet ist, und gib an, wie viele Nullstellen sie hat.}
\begin{geruest}
\gz{a}{$y = (x - 2)^2 + 1$}{\punktfeld \quad \kreuz{oben}\kreuz{unten} Nullstellen: \leerfeld}
\gz{b}{$y = (x - 4)^2 - 3$}{\punktfeld \quad \kreuz{oben}\kreuz{unten} Nullstellen: \leerfeld}
\gz{c}{$y = (x + 1)^2 - 2$}{\punktfeld \quad \kreuz{oben}\kreuz{unten} Nullstellen: \leerfeld}
\gz{d}{$y = -(x - 3)^2 + 4$}{\punktfeld \quad \kreuz{oben}\kreuz{unten} Nullstellen: \leerfeld}
\gz{e}{$y = (x + 2)^2$}{\punktfeld \quad \kreuz{oben}\kreuz{unten} Nullstellen: \leerfeld}
\gz{f}{$y = -(x + 1)^2 - 5$}{\punktfeld \quad \kreuz{oben}\kreuz{unten} Nullstellen: \leerfeld}
\end{geruest}
\end{aufgabe}

\begin{aufgabe}{Ich kann Terme ordnen und zusammenfassen. Fasse zusammen und ordne: zuerst $x^2$, dann $x$, dann die Zahl.}
\begin{teilezwei}
\tz $3x + x^2 + 2 =$ \leerfeld & \tz $5 + x^2 - 4x =$ \leerfeld \\
\tz $x^2 + 2x - 7 + 3x =$ \leerfeld & \tz $2x^2 - 3 + x - x^2 + 8 =$ \leerfeld \\
\tz $4 - 3x - x^2 + 5x =$ \leerfeld & \\
\end{teilezwei}
\end{aufgabe}

\begin{aufgabe}{Ich kann Klammern ausmultiplizieren, binomische Formeln anwenden und $x$ ausklammern. Multipliziere aus und fasse zusammen.}
\begin{teilezwei}
\tz $x \cdot (x + 4) =$ \leerfeld & \tz $(x + 2) \cdot (x + 5) =$ \leerfeld \\
\tz $(x - 3) \cdot (x + 6) =$ \leerfeld & \tz $(x + 5)^2 =$ \leerfeld \\
\tz $(x - 3)^2 =$ \leerfeld & \\
\end{teilezwei}
Klammere $x$ aus.
\begin{teilezwei}
\tz $x^2 + 9x =$ \leerfeld & \tz $x^2 - 2x =$ \leerfeld \\
\end{teilezwei}
\end{aufgabe}

\begin{aufgabe}{Ich finde den Fehler beim Quadrieren einer Klammer. Mia rechnet so:}
\rechnung{(x + 9)^2 &= x^2 + 81}
\begin{teile}
\teil Finde den Fehler, benenne ihn und korrigiere die Rechnung.
\teil Rechne selbst: \ $(x - 8)^2 =$ \leerfeld
\end{teile}
\end{aufgabe}
